\section{后续工作}
\label{sec:future}

在上述八组基线基础上，我们识别出三类后续实验方向，以回应尚未解决的问题。

\subsection{A 组：模型缩放}

\paragraph{A1 —— PicoDet-M。}
PicoDet-M（中等变体，$\approx$4.8\,M 参数，8.2\,MB）预期具有更高精度上限。
我们计划在全量数据集上、采用线性缩放规则超参数进行 300 epoch 双 GPU 训练，估计 mAP 可达 95.5--96.0\%。

\paragraph{A2 —— PP-YOLOE+-S。}
PP-YOLOE+~\citep{xu2022ppyoloe} 是 YOLO 家族的无锚重设计，具有解耦头与任务对齐分配，可缓解 YOLOv3 中观察到的锚框敏感问题。
30 epoch 微调（PaddleDetection 针对 VOC 的调度）可与 YOLOv3-Y4 在相近模型规模下直接对比。

\subsection{B 组：YOLOv3 收敛改进}

\paragraph{B1 —— K-means 锚框聚类。}
默认的 COCO 派生锚框可能不匹配我们实验鼠检测的框宽高比分布。
运行 \texttt{tools/anchor\_cluster.py} 可通过 K-means 得到九个数据集专用锚框，有望加速收敛并提升 Y3/Y4 的上限。

\paragraph{B2 —— Y5：120 epoch 与自定义锚框。}
嵌入 B1 生成的锚框后，120 epoch 训练有望使 YOLOv3 充分收敛（如前所述，Y3/Y4 在 epoch~80 仍在上升），并可能超过 92\% mAP。

\paragraph{B3 —— 框分布分析。}
可视化训练集中边界框尺寸与宽高比的分布，可揭示小目标或遮挡鼠是否存在长尾性能问题。

\subsection{C 组：压缩流程}

\paragraph{C1 —— 量化感知训练（QAT）。}
对 L1 最佳模型使用 PaddleSlim QAT，预期可得约 1.1\,MB 的 INT8 模型，精度损失 $<$0.3\,pp，进一步提升 iOS 推理。

\paragraph{C2 —— YOLOv3 的 FPGM 结构剪枝。}
采用几何中值滤波剪枝~\citep{he2019filter}，在 YOLOv3-Y4 上目标减少 30\% FLOPs 后微调，可将 ONNX 大小从 92\,MB 降至约 60\,MB。
需注意 CoreML 自定义算子问题需单独解决；单靠剪枝无法启用 Metal 加速。

\paragraph{C3 —— 知识蒸馏。}
使用更大教师模型（如 PicoDet-L 或 RT-DETR）向 PicoDet-S 蒸馏，可弥补 L1 与假设更高精度模型之间的差距，而不增加部署成本。

\subsection{推荐优先级}

\begin{enumerate}[leftmargin=1.5em]
  \item B3（框分布分析）—— 零训练成本，立即可获得洞察。
  \item A1（PicoDet-M）—— 预期精度提升最高。
  \item B1+B2（自定义锚框 + Y5）—— 验证数据工程论点。
  \item C1（QAT）—— 进一步提升 iOS 速度的直接路径。
  \item A2、C2、C3 —— 探索性；在现有结果基础上边际回报较低。
\end{enumerate}
